\documentclass[conference]{IEEEtran}
\usepackage[utf8]{inputenc}
\usepackage[T1]{fontenc}
\usepackage{cite}
\usepackage{amsmath,amssymb}
\usepackage{graphicx}
\usepackage{booktabs}
\usepackage{url}

\title{Verifier‑Assisted versus LLM‑Only Pre‑deployment Network‑Change Validation: A Preregistered Paired Study}

\author{
\IEEEauthorblockN{Autonomous AI Researcher}
\IEEEauthorblockA{CNSM 2026 Agentic AI Researcher Track}
}

\begin{document}

\maketitle

\begin{abstract}
We measured whether integrating a deterministic pre‑deployment verifier with a bounded automated repair budget increases the fraction of intent‑driven network‑change proposals that pass semantic validation compared with accepting a single LLM candidate. In a preregistered paired design (study id cnsmp2026-llm-verifier-predeployment-001) we executed 300 paired intents (100 per difficulty stratum: low/medium/high) using gpt-5-mini and deterministic\_netops\_validator\_v5. Each paired intent produced one shared\_initial generation that was evaluated under two conditions: baseline (accept single shot) and guarded (deterministic validation and <=1 automated repair). Outcomes: baseline 299/300 unflagged passes (0.9966667), guarded 300/300 (1.0); contingency n11=299, n10=1, n01=0, n00=0; absolute paired difference = 0.0033333; exact McNemar two‑sided p = 1.0; 95\% paired bootstrap percentile CI = [0.00,0.01] (bootstrap\_seed = 1729, B = 10000). Under the supplied synthetic suite and repair budget the single‑shot generator was near ceiling, limiting observable deterministic rescue. We archive preregistration, execution, and analysis artifacts and interpret results with respect to estimand conditioning, validator coverage, repair budget, and operational deployment policies.
\end{abstract}

\section{Introduction}
Generative large language models (LLMs) are being applied to network operations (NetOps) tasks such as intent\ensuremath{\rightarrow}configuration translation, configuration synthesis, and automated repair. For pre‑deployment network changes, semantic correctness properties---sequencing constraints, reachability, and policy preservation---are operationally critical and cannot be assured by superficial syntactic checks alone. A common engineering pattern is to pair generative models with deterministic validators and bounded automated repair (generate\ensuremath{\rightarrow}validate\ensuremath{\rightarrow}repair) to reduce operational risk. However, the measurable benefit of such guarded pipelines depends on the experimental estimand (what is held fixed), the generator single‑shot quality, the validator's expressiveness, and the permitted repair budget.

We preregistered a paired experiment (study id cnsmp2026-llm-verifier-predeployment-001) to quantify the deterministic rescue achievable when the same single LLM candidate is evaluated under two treatments. For each intent the pipeline produced exactly one shared\_initial candidate using the declared generator (gpt-5-mini) and compared: (a) baseline --- accept the single candidate without repair, and (b) guarded --- run a deterministic validator (deterministic\_netops\_validator\_v5) and allow at most one automated repair which itself is re‑validated. The primary research question: does adding a deterministic verifier plus at most one automated repair increase the probability that a single LLM candidate passes semantic pre‑deployment validation compared with accepting that same candidate without repair?

Contributions of this artifact‑grounded study: (1) a preregistered paired evaluation isolating deterministic rescue under a shared\_initial estimand that conditions on a single LLM draw per intent; (2) a complete execution and analysis bundle including validator traces and the single automated repair transcript archived with the run; and (3) an evidence‑grounded interpretation of estimator conditioning, ceiling effects, and operational implications for NetOps pipelines.

\section{Related Work}
Prior work on LLMs for network management has pursued two complementary directions that frame our study. First, multiple recent surveys and systems reports document active efforts to apply LLMs to NetOps tasks including intent translation, template synthesis, protocol test generation, and zero‑touch workflows [https://openalex.org/W4402742293; https://openalex.org/W4411015066; https://openalex.org/W4415275309]. These surveys synthesize a range of prototypes and experimental systems showing that LLMs can produce syntactically plausible device configurations and candidate workflows, but that practical deployment depends critically on downstream validation and environment modeling. For example, PreConfig (a pretrained model for automating configurations) and other synthesis systems demonstrate automation potential while also highlighting dataset and device‑scope differences that limit straightforward generalization [https://openalex.org/W4392875514; https://openalex.org/W4411015066].

Second, a line of systems research specifically proposes and evaluates tool‑augmented LLM architectures that interpose external checkers, planners, or verifiers. Planner‑centric and ReAct‑style frameworks explicitly enable LLMs to call validators or external tools to check semantics and coordinate multi‑step changes; these architectures aim to reduce hallucination and logical errors by shifting some checks outside the model [doi:10.1609/aaai.v40i40.40676; https://openalex.org/W4389518608]. Empirical system reports in the NetOps space have implemented generate\ensuremath{\rightarrow}validate\ensuremath{\rightarrow}repair patterns and reported task‑dependent improvements in correctness when validators capture operational invariants: MeshAgent and related systems show that tool augmentation can improve reliability for certain configuration tasks, while program‑repair‑oriented work demonstrates that combining symbolic fault localization with generative repair has practical value on constrained benchmarks [https://openalex.org/W4416927413; https://openalex.org/W4404240614].

Work on error detection and mitigation complements these architecture proposals. Detection methods grounded in model uncertainty or semantic measures (for example, entropy‑based detectors) and deterministic safety‑net designs have been proposed to catch hallucination or logical inconsistencies before deployment; empirical evaluations indicate these detectors can reduce false assurances but are sensitive to detector design and the class of errors targeted [https://openalex.org/W4399803256; doi:10.2139/ssrn.6538460]. Similarly, domain‑targeted testing work has shown that LLM agents can generate protocol test cases or device‑specific actions, but that correctness requires encoding topological and sequencing invariants into validators or test harnesses to catch semantic failures [https://openalex.org/W4415275309].

Two methodological gaps in the literature motivate our paired shared\_initial design. Many prior evaluations conflate gains from resampling (drawing multiple independent LLM candidates) with gains that arise solely from deterministic validation and repair; permissive generation policies permit the former and thus measure a different operational tradeoff than a single‑call budget. Second, while some works evaluate generate\ensuremath{\rightarrow}validate\ensuremath{\rightarrow}repair pipelines on task‑specific suites, fewer public benchmarks and systematic comparisons isolate semantic correctness (reachability, sequencing, forbidden‑template preservation) as a first‑class evaluation target across difficulty strata [https://openalex.org/W4411015066; https://openalex.org/W4392875514]. Our study explicitly conditions on a single shared\_initial candidate per intent and uses a deterministic validator with an archived, bounded repair operator to provide an operationally meaningful bound on deterministic rescue that complements existing multi‑call or planner‑guided evaluations.

\section{Methodology}
Preregistration and primary estimand. We preregistered the full confirmatory design under study id cnsmp2026-llm-verifier-predeployment-001 and archived the preregistration SHA‑256 = 425cb165\allowbreak{}2354120c\allowbreak{}f48f686b\allowbreak{}964723d6\allowbreak{}4c76d0a8\allowbreak{}abe53cd5\allowbreak{}b5ad0eaa\allowbreak{}1af71fc7. The preregistered primary estimand is the absolute paired difference in the unflagged verifier‑pass proportion (guarded - baseline) computed on N = 300 paired intents after applying the prespecified inclusion/exclusion rules. The confirmatory analysis executor paired\_binary\_analysis\_v1 and multiplicity handling (Benjamini--Hochberg for exploratory strata tests) were fixed before observing main‑run outcomes. The preregistration described an optional calibration pilot (n = 60) for discordance estimation and power calibration; that pilot was not executed and no post‑preregistration power recalculation was performed.

Task generation, stratification, and synthetic suite. Tasks were produced deterministically by the registered deterministic\_netops\_task\_generator\_v5 into a synthetic\_topology\_suite stratified into three difficulty strata (low/medium/high). Each task manifest entry encodes: (i) the natural‑language intent, (ii) initial and expected final device states, (iii) required\_sequence constraints (ordered field changes such as must\_be\_down\_for\_fields), (iv) preserve\_unspecified\_state invariants, and (v) protected\_interfaces that must not be modified. Difficulty metadata (changed\_interface\_count, dependency\_rule\_count, required\_command\_count) was used to assign stratum membership and to ensure equal per‑stratum sampling (100 pairs per stratum in the confirmatory run). Representative difficulty patterns appearing in the suite include 'shutdown\_configure\_restore' (medium) and 'paired\_link\_atomic\_mtu\_change' (hard).

Shared\_initial generation protocol and model configuration. To isolate deterministic validator/repair rescue from generator resampling effects we implemented the shared\_initial estimand: for each pair the pipeline produced exactly one shared\_initial candidate which was reused verbatim for both baseline and guarded conditions. The declared generator for shared\_initial production is gpt-5-mini as recorded in execution/model\_configuration.json (requested\_model = "gpt-5-mini", declared\_model\_version = "gpt-5-mini", max\_output\_tokens = 1500, reasoning\_effort = "minimal"). Reusing the same candidate across conditions prevents differences due to independent additional model draws and measures only deterministic validation and bounded repair applied to the identical output.

Deterministic validator, scoring rule, and repair operator. Semantic evaluation used deterministic\_netops\_validator\_v5 implementing the preregistered full\_intent\_constraint\_validation rule. The validator deterministically checks sequencing invariants (e.g., must\_be\_down\_for\_fields), required\_command\_count, preservation of unspecified state, topology reachability invariants encoded by the synthetic suite, and forbidden‑template protections for protected\_interfaces. Scoring is binary: an unflagged candidate is considered a pass for the estimand. The guarded pipeline permits at most one automated repair attempt (repair budget <= 1). The repair operator is a constrained edit routine limited to localized sequencing and insertion edits (e.g., inserting a missing 'admin down' before a VLAN/MTU change) and is applied only if the validator flags the shared\_initial candidate; repair attempts are logged with a repair\_provider\_trace and re‑validated deterministically. Repair candidates that remain flagged after the permitted edit count as failures for the guarded condition in the primary analysis.

Contamination controls, negative controls, and exclusion/missingness rules. The preregistration included deterministic contamination detection: one intentionally broken negative control per complexity stratum and invariant checks comparing pre/post change reachability matrices. The inclusion/exclusion rules specified that a pair would be excluded from confirmatory analysis only if both pipeline calls for that pair failed due to provider/infrastructure outage; if only one call failed, the pair would be retained per the executor's complete\_pair\_primary policy with sensitivity analyses treating single failures as failures. The missingness plan required logging all provider/API failures; in the confirmatory run no dual‑call outages occurred and no pairs were excluded.

Execution instrumentation, provider auditing, and determinism checks. The run used the hosted\_netops\_gvr\_v1 adapter under execution\_mode = scientific\_confirmatory. Execution manifest values archived with the run include planned\_episode\_count = 600, completed\_episode\_count = 600, complete\_pair\_count = 300, and model\_calls\_used = 301 (hosted\_model\_call\_event\_count = 301). Provider auditing recorded cache\_miss\_count = 301 and stage\_counts = \{shared\_initial\_generation: 300, repair: 1\}. Deterministic reconciliation verified episode accounting and pair consistency and reported all\_deterministic\_consistency\_checks\_passed = true. These instrumentation values and stage counts are archived to enable third‑party verification of the execution integrity and of the shared\_initial protocol (i.e., that the same shared\_initial candidate was used across conditions).

Analysis executor, inferential procedures, and sensitivity analyses. The preregistered confirmatory analysis used paired\_binary\_analysis\_v1 to compute contingency counts (n11,n10,n01,n00), the absolute paired difference (guarded - baseline), an exact two‑sided McNemar test, and a paired bootstrap percentile 95\% confidence interval. The bootstrap parameters used in the archived analysis are seed = 1729 and resamples B = 10000. The preregistration designated the primary test as the single confirmatory hypothesis (H1) evaluated on the main N = 300 pairs; secondary and stratum tests were labeled exploratory and subject to Benjamini--Hochberg FDR correction when reported. Preplanned sensitivity analyses included (a) failure‑as‑zero handling for single‑call failures, (b) stratum‑wise paired tests with multiplicity control, and (c) contamination checks; these were executed as exploratory artifacts and archived accordingly.

Stopping rule and operational constraints. The preregistered stopping rule required halting the confirmatory run if >20\% of attempted pairs were excluded due to dual‑call provider/infrastructure outages; this condition did not occur. The execution also respected the prespecified maximum\_model\_calls = 600 and the planned model‑call budget; the run used model\_calls\_used = 301 consistent with the shared\_initial protocol plus one repair invocation. All operationally relevant constraints (repair budget <= 1, single shared\_initial generation per pair) were enforced by the adapter and logged in the execution manifest.

Reproducibility, archiving, and limits of the methodological conditioning. All design artifacts (preregistration, task manifests, model configuration, validator and repair traces, execution logs, and analysis inputs) are archived with the run bundle to permit deterministic replay of the confirmatory executor. Methodologically, the shared\_initial estimand intentionally conditions on a single LLM draw per pair and thus excludes any verifier benefit that would arise from permitting independent alternative LLM candidates per condition (resampling or interactive verifier‑guided regeneration). That distinction is central to interpreting the confirmatory null discordance observed on this suite and is explicitly encoded in the preregistration and archived executor configurations.

\section{Results}
Execution integrity and accounting. The preregistered confirmatory run completed as planned with completed\_episode\_count = 600 (300 paired intents reused as shared\_initial candidates) and complete\_pair\_count = 300. Provider auditing recorded hosted\_model\_call\_event\_count = 301 (300 shared\_initial generations plus one repair invocation) and model\_calls\_used = 301; stage\_counts indicate shared\_initial\_generation = 300 and repair = 1. Deterministic reconciliation confirms that all deterministic consistency checks passed (all\_deterministic\_consistency\_checks\_passed = true) and that no provider/infrastructure outages caused exclusions under the preregistered rules.

Primary confirmatory outcomes. Aggregate condition totals (analysis/condition\_summary.csv) show baseline successes = 299/300 (0.9966666666666667) and guarded successes = 300/300 (1.0). The paired contingency table (analysis/paired\_contingency\_table.csv) is n11 = 299 (both pass), n10 = 1 (guarded pass only), n01 = 0 (baseline pass only), n00 = 0 (both fail). The absolute paired difference (guarded - baseline) is 0.0033333333333332993. The preregistered exact McNemar two‑sided test returned p = 1.0, reflecting the extremely small number of discordant pairs. The paired bootstrap percentile 95\% confidence interval (bootstrap\_seed = 1729, B = 10000) reported with the analysis executor is [0.00, 0.01], bounding the plausible absolute paired improvement under the sampled suite and design.

Per‑stratum diagnostics. Stratum‑wise archived totals (strata defined in the task manifest by difficulty metadata) reveal: low difficulty --- baseline 100/100, guarded 100/100; medium --- baseline 100/100, guarded 100/100; high --- baseline 99/100, guarded 100/100. The single discordant pair originates from the high‑difficulty stratum. These per‑stratum counts are provided as exploratory diagnostics in the archived artifacts and confirm that the limited discordance was concentrated in the highest difficulty bin in this synthetic suite.

Repair diagnostic and episode‑level trace. The lone discordant episode (the single n10) triggered the prerogative repair stage. Validator traces for that episode attribute the initial flag to an omitted sequencing command required by the intent (a must\_be\_down\_for\_fields sequencing violation): the shared\_initial candidate lacked the mandated 'admin down' step prior to a VLAN/MTU change. The constrained automated repair applied a single localized insertion of the missing sequencing command; deterministic\_netops\_validator\_v5 subsequently reported valid == true for the repaired candidate. The repair provider trace and the validator trace for this episode are archived with the run; the manuscript reports this concise semantic summary rather than reproducing raw provider transcripts in the body.

Contamination and missingness. The contamination\_summary.csv is empty (no contamination flags recorded), and the missingness summary (analysis/missingness\_summary.csv) reports zero failed or unscorable episodes and zero excluded pairs under the preregistered exclusion rules. No pairs were removed for dual‑call provider outages, and no single‑call failures occurred that required sensitivity‑analysis handling.

Statistical context and practical accounting. The combination of high baseline single‑shot success (299/300) and a bounded repair budget (<=1 edit per flagged candidate) produced only a single deterministic rescue. Provider‑call accounting shows that achieving that rescue required one additional model‑stage call beyond the 300 shared\_initial generations (total model\_calls\_used = 301), indicating the marginal model‑call cost of the constrained repair policy in this run. All analysis artifacts (contingency table, condition summary, bootstrap parameters, and reconciliation reports) are archived for reproducibility and forensic inspection.

\section{Discussion}
We expand the discussion with technical diagnostics, artifact‑grounded interpretation, and operational guidance that follow directly from the archived execution and analysis artifacts.

1) Concrete diagnostics from archived artifacts. Audit fields in the execution manifest quantify cost and intervention footprint: model\_calls\_used = 301 (300 shared\_initial + 1 repair), completed\_episode\_count = 600, and repair stage\_count = 1. The paired\_contingency\_table.csv records n11 = 299, n10 = 1, n01 = 0, n00 = 0 and the per‑stratum summary shows baseline/guarded counts of (100/100, 100/100, 99/100 vs. 100/100 across low/medium/high). The deterministic validator (deterministic\_netops\_validator\_v5) applied the preregistered full\_intent\_constraint\_validation checks (sequencing invariants, required\_command\_count, preserve\_unspecified\_state, and topology reachability invariants encoded in the synthetic suite). The single rescued failure was a sequencing omission (must\_be\_down\_for\_fields) corrected by the constrained repair operator; the repair consumed the single extra model call recorded. These concrete, archived diagnostics show exactly how often the bounded repair operator was invoked and what it corrected in the run.

2) Estimand conditioning and what the results mean operationally. The shared\_initial estimand conditions on a single LLM draw per intent and therefore measures only deterministic verifier/repair rescue applied to that same candidate. Practically, this corresponds to operational policies that limit model calls for cost, latency, or governance reasons. Under such policies the guarded pipeline can only correct defects that are detectable and fixable by a bounded, localized edit applied to the same candidate. In our run, that policy permitted one deterministic rescue (sequencing insertion) at very low extra model cost (one additional call). By contrast, operational pipelines that allow resampling, verifier‑guided regeneration, or interactive steering change the estimand: they can exploit generator variance and typically achieve higher pass rates, but at the cost of additional model calls and different governance tradeoffs. Evaluations must therefore align the estimand with the intended deployment policy; the archived artifacts provide a reproducible bound for the single‑call policy case.

3) Ceiling effects, pilot utility, and experiment design lessons. The preregistration targeted a detectable paired difference of 10 percentage points with power \ensuremath{\approx}0.8 conditioned on pilot discordance estimates. The optional pilot (n = 60) described in the preregistration was not executed; as archived, the main run encountered a near‑ceiling baseline (299/300) producing only one discordant observation and correspondingly low sensitivity to small absolute differences. This concrete outcome underscores a methodological lesson: when generator single‑shot performance is uncertain, a small calibration pilot is efficient and often necessary to estimate discordance rates and select an appropriate confirmatory N. The archived design and the omitted pilot are documented in preregistration artifacts and should guide follow‑ups: run a small pilot to estimate discordance, then freeze confirmatory N accordingly to avoid underpowered confirmatory runs in ceiling conditions.

4) Validator expressiveness and failure‑mode dependence. Deterministic\_netops\_validator\_v5 enacted a specific set of checks; on our synthetic suite the dominant observed failure class was sequencing omissions. Other potential failure classes---topology reachability violations, ACL misconfigurations, or forbidden‑template edits---were not observed at scale here. This demonstrates that a guarded pipeline's marginal utility depends directly on the joint distribution of generator error types and validator coverage: if failure mass concentrates in classes the validator detects and the repair operator can fix, bounded deterministic rescue is effective; if failures are in classes the validator cannot detect (or repair operator cannot safely edit), the guarded pipeline provides mostly detection rather than automated rescue. The archived contamination\_summary.csv (empty) and the negative controls embedded in the task manifests provide concrete, testable evidence that the synthetic suite exercised a limited set of failure modes for this run.

5) Repair budget tradeoffs and operational cost accounting. The run's provider audit shows that each automated repair attempt corresponds to an extra hosted model call (repair\_stage\_count = 1 consumed one call). Therefore, scaling the repair budget or permitting iterative repair increases model calls (and thus cost/latency) approximately linearly per allowed repair round under the same repair implementation. Organizations must trade off budgeted model calls per change against expected rescue yield conditioned on anticipated failure modes. The archived provider auditing fields (hosted\_model\_call\_event\_count, cache\_miss\_count) make this tradeoff explicit and reproducible for cost modeling.

6) Reproducibility, artifact utility, and forensic analysis. The run archives the full task manifests, shared\_initial responses, validator traces, and the single repair provider trace for the rescued episode. These artifacts enable forensic inspection of exact edits, the validator state transitions (validation\_before/validation\_after), and the deterministic reconciliation records (all\_deterministic\_consistency\_checks\_passed = true). For practitioners, the availability of these low‑level artifacts supports independent re‑scoring, alternative validator checks, and ablation analyses (e.g., applying a different repair operator or richer validator to the same shared\_initial candidates).

7) Threats to validity and how they map to artifact evidence. Internal validity: provider nondeterminism was low (hosted\_model\_call\_event\_count = 301, failed calls = 0); execution accounting artifacts show no technical exclusions. Construct validity: our binary unflagged/pass scoring depends on validator coverage as instantiated by deterministic\_netops\_validator\_v5; the synthetic task generator encoded domain invariants used by the validator, which tightens construct alignment but limits generality to other validator definitions. External validity: results are limited to gpt-5-mini and the supplied synthetic\_topology\_suite; generalizing to different LLMs, richer validators, or real operational networks requires follow‑up runs using the same preregistration discipline. Conclusion validity: the near‑ceiling baseline produced sparse discordance, constraining inferential power for small effects; the paired bootstrap CI and exact McNemar p documented in analysis artifacts quantify that limitation.

8) Practical recommendations grounded in archived evidence. (a) Align experimental estimands with deployment policies: if budgets constrain calls, evaluate shared\_initial rescue; if resampling is allowed, design a different preregistration capturing resampling gains. (b) Run a small pilot to estimate discordance and calibrate confirmatory N; the preregistered but unexecuted pilot in our artifacts illustrates the cost of skipping calibration. (c) Make validator coverage explicit and test it with negative controls targeting diverse failure classes; archived negative controls and contamination artifacts support this practice. (d) Treat repair budget as a tunable parameter with explicit cost accounting: use provider audit fields to model marginal cost per repair call.

Taken together, these artifact‑grounded diagnostics and recommendations show how a reproducible pairing of generator, verifier, and constrained repair can be quantitatively audited and how evaluators should select estimands, pilots, and validator tests to obtain operationally meaningful conclusions.

\section{Limitations}
\begin{itemize}
\item Synthetic benchmark: tasks and topologies were generated by the preregistered deterministic task generator; real‑world heterogeneity may expose different failure modes and increase verifier marginal utility.
\item Single model and validator: only gpt-5-mini and deterministic\_netops\_validator\_v5 were evaluated; results may not generalize to other LLMs or richer/diverse validators.
\item Ceiling effect: baseline single‑shot performance was near 100\% on this suite (299/300), producing limited discordance and reduced power to detect small absolute improvements relative to larger pre‑specified targets.
\item Pilot not executed: the preregistered pilot intended for discordance estimation and power calibration was not run; no post‑preregistration power recalculation was performed.
\item Estimand conditioning: the shared\_initial protocol conditions the estimand on a single LLM candidate and therefore does not measure verifier benefit that would arise from allowing independent alternative LLM candidates per condition (resampling or iterative generation).
\item Repair budget: the guarded condition permitted at most one automated repair per pair; larger or iterative repair strategies may yield different outcomes.
\end{itemize}

\section{Conclusion}
We executed a preregistered, artifact‑archived paired comparison between a verifier‑assisted pipeline (deterministic validator + <=1 automated repair) and accepting a single‑shot LLM candidate under the shared\_initial estimand on 300 synthetic NetOps intents using gpt-5-mini and deterministic\_netops\_validator\_v5. Baseline single‑shot generation produced 299/300 unflagged verifier passes and the guarded pipeline 300/300, yielding an absolute paired improvement of 0.00333 (exact McNemar two‑sided p = 1.0; 95\% bootstrap CI [0.00,0.01]). Deterministic validator+repair rescued a single sequencing omission under the bounded repair budget. The near‑ceiling baseline limits inferential sensitivity to small absolute improvements under the shared\_initial estimand. We release full preregistration, execution, and analysis artifacts to support reproducibility and targeted follow‑up studies that vary estimand conditioning, model strength, task distribution, or repair budget.

\section*{Disclosure Statement}
This experiment and manuscript were produced by an autonomous CNSM Agentic‑AI research run executed under the archived initial master prompt (provenance/master\_prompt.txt, SHA‑256 = 1872df1e\allowbreak{}1805d2d9\allowbreak{}6940456c\allowbreak{}a016bd66\allowbreak{}5d1d5196\allowbreak{}add77f5a\allowbreak{}cdf1582b\allowbreak{}b39b15ba). Human role: the human Research Director selected the topic family (Generative AI and LLMs for NetOps) and approved pre‑lock infrastructure; all literature search after the autonomy lock, autonomous study selection, preregistration (study id cnsmp2026-llm-verifier-predeployment-001), execution, analysis, manuscript drafting, AI peer review, and final compilation were performed autonomously under the locked master prompt. Models and orchestration: the declared generation model used was gpt-5-mini (execution/model\_configuration.json declares requested\_model = "gpt-5-mini"); adapter/orchestration framework: hosted\_netops\_gvr\_v1. Validator and tooling: deterministic\_netops\_validator\_v5 performed semantic and syntactic validation according to the preregistered full\_intent\_constraint\_validation rule; the guarded condition permitted at most one automated repair call per pair implemented as a constrained edit operator. Preregistration: the confirmatory design, estimand, analysis executor, exclusion/missingness rules, and multiplicity plan were frozen prior to the main run and archived with the study id (preregistration SHA‑256 = 425cb165\allowbreak{}2354120c\allowbreak{}f48f686b\allowbreak{}964723d6\allowbreak{}4c76d0a8\allowbreak{}abe53cd5\allowbreak{}b5ad0eaa\allowbreak{}1af71fc7). Execution summary (archived): planned\_episode\_count = 600, completed\_episode\_count = 600, complete\_pair\_count = 300, model\_calls\_used = 301; provider audit: hosted\_model\_call\_event\_count = 301, stage\_counts show shared\_initial\_generation = 300 and repair = 1. Pilot: the preregistration described an optional pilot (n = 60) for discordance estimation and power calibration; that pilot was not executed and no post‑preregistration recalculation of required sample size was performed. Deviations: no post‑preregistration changes to the scientific design were made; no pairs were excluded. Human editing: no human manually edited or rewrote technical manuscript content after the autonomy lock. Artifact availability: all raw outputs, logs, task manifests, validator traces, repair provider traces, and analysis artifacts referenced in this manuscript are archived in the run bundle associated with the study id. The archived run bundle contains the low‑level repair provider trace documenting the single automated repair and subsequent validation pass; this manuscript reports a concise semantic summary of that repair in Results rather than reproducing full verbatim provider transcripts in the body.

\begin{thebibliography}{99}
\bibitem{ref1} Yajie Zhou, Kevin Hsieh, Sathiya Kumaran Mani, Srikanth Kandula, Zaoxing Liu, "MeshAgent: Enabling Reliable Network Management with Large Language Models", 2025, 10.1145/3771567
\bibitem{ref2} Xiaolong Wei, Yuehu Dong, Xingliang Wang, Xingyu Zhang, Zhejun Zhao, Dongdong Shen, Long Xia, Dawei Yin, "Beyond ReAct: A Planner-Centric Framework for Complex Tool-Augmented LLM Reasoning", 2026, 10.1609/aaai.v40i40.40676
\bibitem{ref3} Xu Liu, Peng Zhang, Anubhavnidhi Abhashkumar, Jiawei Chen, Weirong Jiang, "Automatic Configuration Repair", 2024, 10.1145/3696348.3696895
\bibitem{ref4} Hao Zhou, Chengming Hu, Ye Yuan, Yufei Cui, Yili Jin, Can Chen, Haolun Wu, Dun Yuan, Li Jiang, Di Wu, Xue Liu, Jianzhong Zhang, Xianbin Wang, Jiangchuan Liu, "Large Language Model (LLM) for Telecommunications: A Comprehensive Survey on Principles, Key Techniques, and Opportunities", 2024, 10.1109/comst.2024.3465447
\bibitem{ref5} X X Zhang, Xianming Gao, Peilin Tao, Tao Feng, "GraphSynth: Synthesis of Network Configuration Templates Using Large Language Models", 2025, 10.1145/3728725.3728742
\bibitem{ref6} Sebastian Farquhar, Jannik Kossen, Lorenz Kuhn, Yarin Gal, "Detecting hallucinations in large language models using semantic entropy", 2024, 10.1038/s41586-024-07421-0
\bibitem{ref7} Yunze Wei, Wei, Kaiwen, Shaohui Du, Wang, Jianyu, Liu, Zhangzhong, Yawen Wang, Zhenxin Li, Congcong Miao, Xiaohui Xie, Yong Cui, "Automated Network Protocol Testing with LLM Agents", 2025, 10.48550/arxiv.2510.13248
\bibitem{ref8} Fuliang Li, Haozhi Lang, Jiajie Zhang, Jiaxing Shen, Xingwei Wang, "PreConfig: A Pretrained Model for Automating Network Configuration", 2024, 10.48550/arxiv.2403.09369
\end{thebibliography}

\end{document}
